
\subsection{Upper-limb Biomechanical Models}

  All four models share a single base kinematic structure, the 36-DOF
  ISB-convention whole-body model. Each model is obtained by applying a set of
  targeted, literature-anchored modifications to the modern-human reference while leaving the joint topology unchanged. Edits fall
  into four families: (i)~\emph{segment geometry} (link/joint origins encoding
  bone lengths and orientations), (ii)~\emph{mass--inertia} properties,
  (iii)~\emph{joint orientation offsets} (RPY terms encoding humeral torsion and
  cranial shoulder tilt), and (iv)~\emph{actuation limits} (joint effort and
  velocity bounds used as a proxy for muscular capacity). Modifications are
  applied bilaterally and symmetrically to the left and right limbs.
  Table~\ref{tab:morph} summarises the resulting parameter set.




  \begin{table}[h]
  \centering
  \caption{Implemented upper-limb parameters across the four models, relative to
  the modern-human ISB reference. Effort in N\,m, velocity in rad\,s$^{-1}$.
  Non-human shoulder/elbow efforts were rescaled to comparative muscle-force
  data (see Table~\ref{tab:rom}); the modern-human column is the unchanged ISB
  source.}
  \label{tab:morph}
  \small
  \begin{tabular}{@{}lcccc@{}}
  \toprule
  Parameter & Modern human & \textit{H.\,naledi} & Neanderthal & Chimpanzee \\
  \midrule
  \multicolumn{5}{@{}l}{\textit{Segment geometry}}\\
  Humeral length (m)            & 0.276 & 0.256 & 0.248 & 0.283 \\
  Humeral torsion offset (rad)  & 0.00  & 0.925 & 0.00  & 0.08  \\
  Shoulder cranial tilt (rad)   & 0.00  & $-0.366$ & 0.00 & $-0.25$ \\
  Clavicle origin, lateral (m)  & 0.176 & 0.160 & 0.210 & 0.180 \\
  \midrule
  \multicolumn{5}{@{}l}{\textit{Mass--inertia}}\\
  Thorax mass (kg)              & 11.97 & 9.50  & 15.50 & 11.00 \\
  Upper-arm mass (kg)           & 1.80  & 1.42  & 1.80  & 1.60  \\
  Forearm mass (kg)             & 1.28  & 1.28  & 1.66  & 1.28  \\
  \midrule
  \multicolumn{5}{@{}l}{\textit{Actuation limits (effort / velocity)}}\\
  Thoracic flexion              & 190 / 5.2 & 120 / 6.0 & 250 / 4.0 & 180 / 8.0 \\
  Clavicle                      & 100 / 14.0 & 80 / 12.0 & 150 / 10.0 & 120 / 20.0 \\
  Shoulder ($X$, abduction)     & 71 / 6.0  & 71 / 6.0  & 92 / 5.0  & 71 / 6.0  \\
  Shoulder ($Y$, elevation)     & 52 / 6.0  & 44 / 6.0  & 52 / 6.0  & 73 / 15.0 \\
  Elbow flexion                 & 77 / 19.9 & 70 / 19.9 & 115 / 12.0 & 108 / 25.0 \\
  \bottomrule
  \end{tabular}
  \end{table}

    \subsubsection*{Rationale for the four models}
  We compare four upper-limb morphologies chosen to bracket the range of
  variation relevant to manual tool use, spanning from a living, habitual
  tool-maker to a living ape outgroup, with two extinct hominins of contrasting
  build in between (Fig.~\ref{fig:species}):
  \begin{itemize}
    \item \textbf{Modern human} (\textit{Homo sapiens}) --- the reference
    morphology and the only taxon for which in-task demonstration data are
    available; a habitual, dexterous maker and user of flaked stone tools.
    \item \textbf{Neanderthal} (\textit{H.\ neanderthalensis}) --- a closely
    related, robustly built hominin and an established maker of
    Mousterian/Levallois lithic industries. It tests whether a
    high-robusticity, high-strength regime shifts the optimal manipulation
    strategy.
    \item \textbf{\textit{Homo naledi}} --- a small-bodied hominin with a
    \emph{mosaic} upper limb: human-like in some hand and wrist features yet
    primitive in others (a low-torsion humerus, curved phalanges, and a
    cranially oriented, ``shrugged'' shoulder). Its tool-related behaviour
    remains debated, making it an informative intermediate between the human
    and ape conditions.
    \item \textbf{Chimpanzee} (\textit{Pan troglodytes}) --- the closest living
    relative and a tool \emph{user} (e.g.\ stone hammers for nut-cracking) but
    not a knapper of flaked stone tools. It serves as the extant non-hominin
    outgroup and a control for what an ape-grade upper limb predicts.
  \end{itemize}
  Together these taxa span a graded axis of upper-limb form---in segment
  proportions, robusticity, joint orientation, and range of motion---while
  performing the identical contact-rich shaving task. This lets us ask whether
  morphology alone shifts the cost function recovered by inverse reinforcement
  learning, and whether any such shift is consistent with the archaeological and
  behavioural record of who did, and did not, manufacture flaked stone tools.
    \begin{figure}[h]
    \centering
     \includegraphics[width=0.23\textwidth]{figures/models/neanderthalensis_JG_Recon_Head_CC_3qtr_lt_l.jpg.png}
    \includegraphics[width=0.23\textwidth]{figures/models/H naledi reconstruction IMG_2355a l.jpg.png}
    \includegraphics[width=0.23\textwidth]{figures/models/f369758da04e040151acc19bf7a552cb.jpg}

    \caption{The models compared in this study (left to right):
    Neanderthal, \textit{Homo naledi}, and chimpanzee. Source: https://humanorigins.si.edu/evidence/human-fossils/species , https://www.thegreatprojects.com/blog/fascinating-chimpanzees-facts]}
    \label{fig:species}
  \end{figure}

  \newcommand{\armpanel}[9]{%
  % xscale=-1 mirrors the limb to the subject's RIGHT arm (the analysed chain);
  % node text is unaffected by the transform, so labels stay upright.
  \begin{scope}[shift={(#1,0)},xscale=-1]
    % --- torso silhouette (tapered; width encodes thoracic mass) ---
    \coordinate (BL) at (-#9/2,0);        \coordinate (BR) at (#9/2,0);
    \coordinate (TR) at (#9*0.31,3.2);    \coordinate (TL) at (-#9*0.31,3.2);
    \draw[#8,fill=#8!10,rounded corners=3pt,line width=0.6pt]
          (BL) -- (BR) -- (TR) -- (TL) -- cycle;
    \draw[#8,line width=0.6pt] (-0.11,3.2) -- (-0.09,3.5)
                               ( 0.11,3.2) -- ( 0.09,3.5);   % neck
    \draw[#8,fill=#8!10,line width=0.6pt] (0,3.82) circle (0.32); % head
    % --- shoulder (raised by cranial tilt) ---
    \coordinate (S) at (#9*0.31, {3.2 + #4*0.012});
    \pgfmathsetmacro{\ha}{-90 + #4}     % humerus direction
    \pgfmathsetmacro{\fa}{-50 + #4}     % forearm direction (compact flexion)
    \coordinate (E) at ($(S)+(\ha:#3)$);
    \coordinate (H) at ($(E)+(\fa:2.1)$);
    % --- limb segments (width encodes robusticity) ---
    \draw[#8,line width=#6 pt,line cap=round] (S) -- (E);
    \draw[#8,line width=#7 pt,line cap=round] (E) -- (H);
    \fill[#8] (S) circle (2.6pt);  \fill[#8] (E) circle (2.2pt);
    % --- stick held at the hand ---
    \draw[brown!70!black,line width=2.2pt,line cap=round]
          ($(H)+(190:0.4)$) -- ($(H)+(-10:0.4)$);
    % --- humeral length label ---
    \pgfmathsetmacro{\lm}{#3/10}
    \node[#8,font=\scriptsize,left=1pt] at ($(S)!0.5!(E)$) {\pgfmathprintnumber[precision=3]{\lm}\,m};
    % --- humeral torsion: drawn ON the humerus shaft (not the elbow) and
    %     spelled out, so the value is not misread as the elbow flexion angle ---
    \coordinate (Tpos) at ($(S)!0.6!(E)$);
    \ifnum#5>10
       \draw[#8,->,line width=0.7pt]
          ([shift=(40:0.24)]Tpos) arc[start angle=40,end angle=300,radius=0.24];
       \node[#8,font=\scriptsize,right=6pt] at (Tpos) {torsion #5\textdegree};
    \else
       \node[black!55,font=\scriptsize,right=6pt] at (Tpos) {torsion #5\textdegree};
    \fi
    % --- cranial tilt: beside the shoulder joint, clear of the head ---
    \ifnum#4>0 \node[#8,font=\scriptsize] at ($(S)+(0,0.28)$) {tilt #4\textdegree}; \fi
    % --- taxon label ---
    \node[#8,font=\bfseries\small] at (0,-0.8) {#2};
  \end{scope}}



  \begin{figure}[t]
  \centering
  \resizebox{\textwidth}{!}{%
  \begin{tikzpicture}[>=stealth]
    %          x   name                 hum  tilt tor  upW  foW  colour      torsoW
    \armpanel{0}  {Modern human}       {2.76}{0} {0}  {2.4}{1.8}{cbModern} {1.40}
    \armpanel{4}  {\textit{H.\,naledi}}{2.56}{21}{53} {1.5}{1.5}{cbNaledi} {1.10}
    \armpanel{8}  {Neanderthal}        {2.48}{0} {0}  {2.6}{3.0}{cbNeander}{1.85}
    \armpanel{12} {Chimpanzee}         {2.83}{14}{5}  {2.0}{1.8}{cbChimp}  {1.30}
  \end{tikzpicture}}
  \caption{Comparative upper-limb kinematics of the four models, drawn in a
  common sharpening pose. All share the 36-DOF ISB base topology; panels differ
  only in the literature-anchored parameters of Table~\ref{tab:morph}.
  \emph{Visual encoding:} humeral segment drawn to scale and labelled (m);
  limb line-thickness $\propto$ segment robusticity (mass); torso width
  $\propto$ thoracic mass; the curved arrow ($\circlearrowleft$, degrees) marks
  humeral torsion (an axial, out-of-plane rotation); a raised shoulder denotes
  cranial tilt; the brown bar is the worked stick. \textit{H.\,naledi} is
  distinguished by its $53^\circ$ torsion offset and shrugged shoulder, the
  Neanderthal by its short, robust forearm and heavy thorax, the chimpanzee by
  its long humerus and high suspensory shoulder.}
  \label{fig:morph}
  \end{figure}

  \paragraph{\textit{Homo naledi}: the climber}
  Parameters for \textit{H.\,naledi} are derived from specimen U.W.\,101-283 as
  reported by \citet{feuerriegel2017}. The humeral length is set to
  $256$\,mm (Table~3 of \citealp{feuerriegel2017}), substantially shorter than
  the modern human mean ($\sim$\,323\,mm), shortening the effective lever and
  requiring the model to operate closer to the workpiece. The defining
  modification is a low humeral torsion of $91.0^\circ$ versus $144.1^\circ$ in
  modern humans (Table~4), encoded as a $\sim$$53^\circ$ internal-rotation offset
  ($\mathrm{rpy}_z = 0.925$). This rotates the elbow medially (``pigeon-toed''),
  so that maintaining a tool parallel to the chest is no longer kinematically
  natural. A superiorly and laterally positioned scapula
  (glenoid angle $121.1^\circ$ vs.\ $142.4^\circ$; Table~8) is represented by a
  $\sim$$21^\circ$ cranial shoulder tilt, producing a ``shrugged'' posture.
  Reduced robusticity (robusticity index $0.1835$ vs.\ $0.2055$; Table~5) is
  reflected in lower segment masses and effort limits. We therefore predict the
  inverse-reinforcement-learning (IRL) recovered reward to favour
  \emph{high-frequency, low-pressure} strokes and elevated, near-overhead working
  heights, consistent with a limb adapted for suspension rather than power
  manipulation.

\paragraph{Neanderthal: the power-scraper}
  Neanderthal geometry follows \citet{degroote2011} and the activity-based
  muscular argument of \citet{shaw2012}. Distal shortening (ulna $247.91$\,mm,
  radius $227.76$\,mm; \citealp[Table~4]{degroote2011}) is a cold-climate
  adaptation producing a shorter but stiffer lever favourable for high-pressure
  forces applied close to the body. Forearm mass is increased by $\sim$30\%
  (``hyper-polar robusticity''; \citealp[p.~397]{degroote2011}) with the
  centre of mass shifted laterally to approximate the bowed, laterally curved
  radius, which increases rotational moment arms. The barrel-shaped thorax is
  modelled as a heavy ($15.5$\,kg), stiff anchor (thoracic effort $250$\,N\,m
  over a restricted $\pm0.6$\,rad range), and a $\sim$15\% longer clavicle
  displaces the shoulder laterally. Elbow ($+50\%$) and shoulder-abduction
  ($71{\to}92$\,N\,m, $\sim$$1.3\times$ human) effort limits encode the elevated pectoralis-major and
  anterior-deltoid demand reported during scraping
  (\citealp[Tables~1--2]{shaw2012}). We predict an IRL reward dominated by
  \emph{force consistency} over peak speed, with coordinated torso-rotation and
  arm-extension---a whole-body ``cleaving'' strategy.

  \paragraph{Chimpanzee: the explosive climber control}
  \textit{Pan} parameters use the comparative control data of
  \citet{feuerriegel2017} and primate muscle architecture from
  \citet{oneill2017}. The humerus is the longest of the set ($283.5$\,mm;
  \citealp[Table~3]{feuerriegel2017}), increasing whole-limb inertia and
  penalising the short, precise strokes that whittling requires. Humeral torsion
  is near-modern ($139.0^\circ$; Table~4), so the elbow orientation
  ($\mathrm{rpy}_z = 0.08$) resembles ours; the chimpanzee's predicted task
  difficulty therefore arises not from the arm but from hand/wrist alignment
  constraints. The shoulder is highly cranially oriented and powerful, but its
  strength advantage is calibrated to \citet{oneill2017}, who found chimpanzee
  muscle exceeds human by only $\approx$$1.35$--$1.5\times$ in maximum
  \emph{dynamic} force/power (driven by a $67\%$ vs.\ $40\%$ fast-twitch fibre
  fraction and longer fascicles), \emph{not} by an exceptional static torque.
  We therefore set the shoulder-elevation effort to $73$\,N\,m
  ($\sim$$1.4\times$ human) and carry the dynamic edge through the elevated
  velocity limit ($15$\,rad\,s$^{-1}$), rather than the previously assumed
  $2$--$3\times$ static torque. The chimpanzee's manipulation deficit is encoded
  \emph{distally} instead: a restricted wrist (dorsiflexion capped at
  $\sim$$20^\circ$) and reduced forearm pronation--supination
  (Table~\ref{tab:rom}) reproduce the hand/wrist-alignment constraint. This
  replaces the earlier elbow ``bony stop,'' which \citet{straus1940} (via
  \citealp{knuckle2020}) showed to be a soft-tissue, not osseous, limit. We
  expect the IRL reward to expose a mismatch: explosive dynamic capacity with a
  high energetic penalty for sustained, repetitive manipulation, yielding
  ``hooking'' rather than ``whittling'' trajectories.

  \begin{figure}
      \centering
      \includegraphics[width=0.95\linewidth]{figures/models/reach2.png}
      \caption{Reach of different models}
      \label{fig:placeholder}
  \end{figure}

  \begin{table}[h]
  \centering
  \caption{Forward characterization of the four models, computed
  directly from the URDFs (right limb) rather than from demonstrations.
  \emph{Max reach}: farthest right-hand position from the shoulder over $6000$
  arm configurations sampled uniformly within the joint limits. \emph{Tool-axis
  cone}: angular spread of the achievable hand long-axis (a dexterity proxy from
  the mean resultant length of the sampled orientations). \emph{Static effort}
  $\lVert\tau\rVert$: gravity-compensation torque norm holding a fixed
  mid-reach target ($[0.35,0,1.2]$\,m). \emph{Capacity used}
  $\sum|\tau_i|/\tau_{i,\max}$: summed per-joint fraction of the effort limits at
  that pose. Reach and effort separate the models by segment geometry and limb
  mass; the contracted tool-axis cones of the chimpanzee and \textit{H.\,naledi}
  reflect their reduced distal range of motion; and the chimpanzee's low
  capacity-used value follows from the muscle-force-based effort rescaling of
  Table~\ref{tab:rom}.}
  \label{tab:capabilities}
  \small
  \begin{tabular}{@{}lcccc@{}}
  \toprule
  Model & Max reach & Tool-axis cone & Static effort & Capacity used \\
        & (cm)      & (deg)          & $\lVert\tau\rVert$ (N\,m) & $\sum|\tau_i|/\tau_{i,\max}$ \\
  \midrule
  Modern human        & 79 & 163 & 8.9 & 0.22 \\
  \textit{H.\,naledi} & 58 & 139 & 6.3 & 0.21 \\
  Neanderthal         & 76 & 159 & 9.7 & 0.21 \\
  Chimpanzee          & 62 & 148 & 5.7 & 0.15 \\
  \bottomrule
  \end{tabular}
  \end{table}


  \subsubsection{Joint range-of-motion and actuation-limit calibration}
  \label{sec:rom}
  The base ISB joint \emph{ranges} are generous and identical across taxa, so
  the original models differentiated species almost entirely through segment
  geometry and inertia. To let the reachable workspace and manipulation
  capability also reflect morphology, we calibrated the distal-joint range
  limits and rescaled the actuation efforts against measured comparative data
  (Table~\ref{tab:rom}). The modern-human model is left at its original ISB
  source values and shown only as the reference baseline. Three points warrant
  emphasis for verification:
  \begin{enumerate}
    \item \emph{Effort rescaling.} The chimpanzee ``super-strength'' advantage
    of \citet{oneill2017} is $\approx$$1.35$--$1.5\times$ in maximum
    \emph{dynamic} force/power, \emph{not} in static torque. We scale the
    chimpanzee shoulder/elbow efforts to $\sim$$1.4\times$ human (shoulder
    elevation $52{\to}73$\,N\,m; elbow $77{\to}108$\,N\,m) and carry the dynamic
    edge through velocity limits. Neanderthal and \textit{H.\,naledi} shoulder
    efforts are scaled to $\sim$$1.3\times$ and $\sim$$0.85\times$ human.
    \item \emph{Wrist as the chimpanzee bottleneck.} The knuckle-walking
    chimpanzee wrist has a measured dorsiflexion range of only
    $\sim$$5$--$20^\circ$ \citep{knuckle2020}; we cap \texttt{wrist\_Z}
    extension at $0.35$\,rad ($20^\circ$) and reduce forearm rotation
    (\texttt{elbow\_Y}) and wrist deviation (\texttt{wrist\_X}). \citet{straus1940} showed the limit is
    soft-tissue, not osseous.
    \item \emph{Inference and sign conventions.} Fossil ROM cannot be measured;
    \textit{H.\,naledi} and Neanderthal ranges are set near-human and flagged as
    inferred. Because the left/right \texttt{wrist\_X} and \texttt{elbow\_Y}
    axes are mirrored, \texttt{wrist\_X} is kept symmetric; the asymmetric
    \texttt{wrist\_Z} cap is consistent only because both wrists share a $+z$
    axis. The assumption that the \texttt{wrist\_Z} upper bound is
    extension/dorsiflexion should be confirmed in the kinematic viewer.
  \end{enumerate}

  \begin{table}[h]
  \centering
  \caption{Calibrated distal-joint range-of-motion and effort limits for the
  three non-human models. Ranges as $[\text{lower},\text{upper}]$ in radians
  (degrees in parentheses); efforts in N\,m. The modern-human column is the
  unchanged ISB reference. Confidence: \textbf{M}\,=\,measured comparative data,
  \textbf{I}\,=\,inferred from morphology.}
  \label{tab:rom}
  \small
  \begin{tabular}{@{}llcccc@{}}
  \toprule
  Joint (DOF) & & Human (ref.) & \textit{H.\,naledi} & Neanderthal & Chimpanzee \\
  \midrule
  Wrist flex/ext   & rad   & $[-1.57,1.57]$ & $[-1.13,1.22]$ & $[-1.13,1.22]$ & $[-1.05,0.35]$ \\
                                      & conf. & ---            & I              & I              & \textbf{M} \\
  Wrist dev.\      & rad   & $\pm0.785$     & $\pm0.45$      & $\pm0.45$      & $\pm0.26$ \\
                                      & conf. & ---            & I              & I              & I \\
  Forearm pron/sup, upper) & rad & $3.14$   & $2.62$         & $2.79$         & $1.40$ \\
                                      & conf. & ---            & I              & I              & I \\
  Elbow ext.\ floor , lower) & rad & $0$     & $0$            & $0$            & $0$ \\
  \midrule
  Shoulder elevation ) & N\,m & 52 & 44 & 52 & 73 \\
  Shoulder abduction ) & N\,m & 71 & 71 & 92 & 71 \\
  Elbow flexion          & N\,m & 77 & 70 & 115 & 108 \\
  \bottomrule
  \end{tabular}
  \end{table}

  \subsubsection{Implementation notes and limitations}
  (1)~Effort limits are partly anchored to comparative muscle-force data
  (\citealp{oneill2017}; Table~\ref{tab:rom}) but remain joint-level
  \emph{proxies} for muscular capacity, not measured joint torques. (2)~The
  current URDFs are bilaterally symmetric; the right-dominant strength asymmetry
  (24--57\%) reported for Neanderthals by \citet{shaw2012} is \emph{not} yet
  implemented and is left for future work. (3)~In the Neanderthal model the
  humeral (upper-arm) link length is set to the published \emph{ulna} value
  ($248$\,mm); this conflates forearm and upper-arm shortening and should be
  revisited if segment-specific lever effects are analysed. (4)~Fossil
  (\textit{H.\,naledi}, Neanderthal) range-of-motion limits are inferred rather
  than measured and set near-human; only the chimpanzee wrist-extension cap
  rests on measured kinematic data \citep{knuckle2020}.

% ===========================================================================
% REFERENCES TO ADD TO THE .bib AND DOUBLE-CHECK BEFORE SUBMISSION
% [confirmed] = citation verified;  [VERIFY] = confirm author/year/page.
% ===========================================================================
% oneill2017  [confirmed]
%   O'Neill MC, Umberger BR, Holowka NB, Larson SG, Reiser PJ (2017).
%   "Chimpanzee super strength and human skeletal muscle evolution."
%   PNAS 114(28):7343-7348.  https://www.pnas.org/doi/10.1073/pnas.1619071114
%   -> chimp muscle ~1.35x human max dynamic force/power; fast-twitch 67% vs 40%.
%
% knuckle2020  [VERIFY authors]
%   "The biomechanics of knuckle-walking: 3-D kinematics of the chimpanzee and
%   macaque wrist, hand and fingers." J. Exp. Biol. 223(14):jeb224360 (2020).
%   https://journals.biologists.com/jeb/article/223/14/jeb224360
%   -> chimp wrist max extension (dorsiflexion) only ~5-20 deg.
%
% straus1940  [VERIFY full ref]
%   Straus WL Jr (1940). Passive chimpanzee wrist extension increases as digital
%   flexors are sequentially transected -> extension is soft-tissue limited, not
%   osseous. (cited via knuckle2020).
%
% wristmoments2024  [VERIFY]
%   "A cadaveric study of wrist-joint moments in chimpanzees and orangutans with
%   implications for the evolution of knuckle-walking." J. Hum. Evol. (2024).
%   https://www.sciencedirect.com/science/article/pii/S0047248424001088
%
% myatt2012  [VERIFY]
%   Myatt JP et al. (2012). "Functional adaptations in the forelimb muscles of
%   non-human great apes." J. Anat.
%   https://onlinelibrary.wiley.com/doi/full/10.1111/j.1469-7580.2011.01443.x
%   -> longer ape pronator fascicles (force over greater pron/sup range).
%
% norkinwhite  [confirmed standard ref] -- HUMAN normative ROM (model UNCHANGED;
%   listed only as the baseline reference column)
%   Norkin CC, White DJ. "Measurement of Joint Motion: A Guide to Goniometry."
%   F.A. Davis.  -> wrist flex ~73-80, ext ~70, radial dev ~20, ulnar dev ~30;
%   forearm pron/sup ~80/80 deg; shoulder flex/abd 180, ext rot 90 (AAOS).
%
% Human peak isometric torque (context for the effort baselines; not applied to
% the human model, which keeps its original ISB efforts):
% wristtorque1996  [VERIFY author/year]  PubMed 8884484 -- wrist flex ~12 Nm, ext ~7 Nm.
% elbowtorque2018  [VERIFY] systematic review -- elbow flexion ~46 Nm at 90 deg.
% shouldertorque2025 [VERIFY] J. Funct. Morphol. Kinesiol. (PMC12821543)
%   -- shoulder flexion ~77 Nm at 90 deg.
% ===========================================================================